\begin{homeworkProblem}[4]
  Sejam $V$ um espaço vetorial normado e $\phi$ um funcional linear em $V$.
  \begin{enumerate}[a)]
    \item Mostre que as quantidades
      \begin{align*}
        \sup\left\{ |\phi(x)|:x\in V,\norm{x}\leq 1 \right\},\quad\sup\left\{ |\phi(x)|:x\in V,\norm{x}=1 \right\},\quad\sup_{x\neq 0}\frac{|\phi(x)|}{\norm{x}},
      \end{align*}
      são iguais e que são finitas se, e somente se, $\phi$ é limitado.
    \item Mostre que o conjunto dos funcionais lineares limitados em $V$ é um espaçõ métrico completo. 
  \end{enumerate}
  \begin{solution}
  \begin{enumerate}[a)]
    \item Note que se denotamos $y=\frac{x}{\norm{x}}$, temos que como $\phi$ é um funcional linear, então 
      \begin{align*}
        \frac{|\phi(x)|}{\norm{x}}=\left| \phi\left( \frac{x}{\norm{x}} \right) \right|=\left| \phi(y) \right|.
      \end{align*}
      Logo temos que como $\norm{y}=1$ para todo $x\in V\setminus\{0\}$, então podemos garantir que
      \begin{align*}
        \sup_{x\neq 0}\frac{|\phi(x)|}{\norm{x}}=\sup_{\norm{x}=1}|\phi(x)|.
      \end{align*}
      Agora, note que como $\{x\in V:\norm{x}=1\}\subseteq\left\{ x\in V:\norm{x}\leq 1 \right\}$, então temos que
      \begin{align*}
        \sup\left\{ |\phi(x)|:x\in V,\norm{x}= 1 \right\}\leq        \sup\left\{ |\phi(x)|:x\in V,\norm{x}\leq 1 \right\}.
      \end{align*}
      Por outro lado, note que se $\sup\left\{ |\phi(x)|:x\in V,\norm{x}\leq 1 \right\}=0$, então temos que $|\phi(x)|=0$ para todo $x\in V$, pois $\phi(0)=0$ e para todo $x\in V\setminus\{0\}$ temos que se $y=\frac{x}{\norm{x}}$, então $\phi(x)=\phi(y)\norm{x}=0$, pelo que a igualdade com $\sup\left\{ |\phi(x)|:x\in V,\norm{x}=1 \right\}$ é trivial. Portanto, vamos considerar que $\sup\left\{ |\phi(x)|:x\in V,\norm{x}\leq 1 \right\}\neq 0$.\\
      Segue-se do anterior que
      \begin{align*}
        \sup\left\{ |\phi(x)|:x\in V,\norm{x}\leq 1 \right\}=        \sup\left\{ |\phi(x)|:x\in V\setminus\{0\},\norm{x}\leq 1 \right\}.
      \end{align*}
      Portanto, note que se $x\in V\setminus\{0\}$ com $\norm{x}\leq 1$ se cumpre que
      \begin{align*}
        |\phi(x)|\leq \frac{|\phi(x)|}{\norm{x}}.
      \end{align*}
      Logo como $\{x\in V:\norm{x}\leq 1\}\subseteq V\setminus\{0\}$, temos que 
      \begin{align*}
        \sup\left\{ |\phi(x)|:x\in V\setminus\{0\},\norm{x}\leq 1 \right\}\leq\sup_{x\neq 0}\frac{|\phi(x)|}{\norm{x}}.
      \end{align*}
      Pelo que podemos concluir que
      \begin{align*}
        \sup\left\{ |\phi(x)|:x\in V,\norm{x}\leq 1 \right\}=\sup_{x\neq 0}\frac{|\phi(x)|}{\norm{x}}.
      \end{align*}
      O que nos permite garantir que as 3 quantidades são iguais.\\
      Para finalizar, note que nos definimos $\left\{ \phi \right\}=\sup\left\{ |\phi(x)|:x\in V,\norm{x}\leq 1 \right\}$, pelo que podemos afirmar que $\phi$ é limitado se, e somente se, as 3 quantidades são finitas.
    \item Seja $L(V)$ o conjunto dos funcionais lineares limitados com a norma
      \begin{align*}
        \norm{\phi}=\sup_{\substack{x\in V\\\norm{x}=1}}|\phi(x)|.
      \end{align*}
      Suponha $\{\phi_n\}_{n\in\mathbb{Z}^{+}}\subseteq L(V)$ sequência de Cauchy. Vamos ver que $\phi_n\to\phi\in L(V)$.\\
      Seja $x\in V\setminus\{0\}$. Note que como $\{\phi_n\}_{n\in\mathbb{Z}^{+}}$ é de Cauchy, temos que dado $\epsilon>0$ existe $N>0$ tal que se $n,m\geq N$, então
      \begin{align*}
        \norm{\phi_n-\phi_m}<\frac{\epsilon}{\norm{x}_{V}}.
      \end{align*}
      Portanto, temos que  
      \begin{align*}
        \left|\phi_n(x)-\phi_m(x)\right|&\leq \norm{\phi_n-\phi_m}\norm{x}_{V}\\
        &<\frac{\epsilon}{\norm{x}_{V}}\norm{x}_{V}\\
        &<\epsilon.
      \end{align*}
      Logo $\{\phi_n(x)\}_{n\in\mathbb{Z}^{+}}\subseteq \mathbb{R}$ é uma sequência de Cauchy, mas como $\mathbb{R}$ é completo, temos  que existe $y_x\in \mathbb{R}$ tal que $\phi_{n}(x)\to y_x$.\\
      Assim, defina $\phi:V\to\mathbb{R}$ dado por $\phi(x)=y_x$.\\
      Agora, vamos ver a convergência em norma.\\
      Como $\{\phi_n\}_{n\in\mathbb{Z}^{+}}$ é de Cauchy temos que dado $\epsilon>0$, existe $N$ tal que para todo $n,m\geq N$,
      \begin{align*}
        \norm{\phi_m - \phi_n} < \epsilon.
      \end{align*}
      Em particular, fixando $n$ y pegando límite quando $m\to \infty$, temos que para todo $x\in V$ tal que $\norm{x}=1$ se cumpre que
      \begin{align*}
        \left| \phi(x)-\phi_n(x) \right|=\lim_{m \to \infty}\left| \phi_m(x)-\phi_n(x) \right|\leq \lim_{m \to \infty}\norm{\phi_m-\phi_n}\norm{x}_{V}< \epsilon.
      \end{align*}
      Logo como $\epsilon$ é independente de $x$, temos que pegando $\sup$ sobre $x$ con $\norm{x}_V = 1$ obtemos que
      \begin{align*}
        \norm{\phi-\phi_n}<\epsilon,
      \end{align*}
      para todo $n\geq N$. Portanto
      \begin{align*}
        \lim_{\epsilon \to 0}\norm{\phi_n-\phi}=0.
      \end{align*}
      Logo usando o anterior temos que dado $\epsilon>0$ existe $N\in\mathbb{Z}^{+}$ tal que se $n\geq N$, então
      \begin{align*}
        \norm{\phi-\phi_n}<\epsilon.
      \end{align*}
      Logo usando a desigualdade triangular temos que como $\phi_n$ é um funcional linear limitado, então 
      \begin{align*}
        \norm{\phi}\leq \norm{\phi-\phi_n}+\norm{\phi_n}<\epsilon+\norm{\phi_n}<+\infty.
      \end{align*}
      Depois, note que $\phi$ é linear, pois dados $x,y\in V$ e $c$ escalar temos que
      \begin{align*}
        \phi(x+cy)=\lim_{n \to \infty}\phi_n(x+cy)=\lim_{n \to \infty}\phi_n(x)+c\phi_n(y)=\phi(x)+c\phi(y).
      \end{align*}
      Pelo que temos garantido que $\phi\in L(V)$ e que $\phi_n\to\phi$ em $L(V)$ quando $n\to\infty$ o que nos permite concluir que $L(V)$ é um espaço métrico completo. 
  \end{enumerate} 
  \end{solution}
\end{homeworkProblem}
